% Carta de presentación en español. Así como está sirve para cualquier empresa.
% Para personalizarla, escribe la empresa y el puesto entre las llaves; si los
% dejas vacíos, la carta se mantiene genérica.
% Compilar: powershell -ExecutionPolicy Bypass -File latex\compilar.ps1 -Documentos carta-es
\documentclass[es,11pt]{yaac-rl}

\empresa{}   % p. ej. \empresa{Contoso}
\puesto{}    % p. ej. \puesto{Desarrollador .NET Senior}

\geometry{margin=2cm}
\hypersetup{
  pdftitle  = {Rafael León Torres: carta de presentación},
  pdfauthor = {Rafael León Torres},
  pdflang   = es-MX
}

% !TeX root = ../cv-es.tex
% Encabezado en español: lo usan el CV y la carta (cv-es.tex y carta-es.tex).
\name{}{Rafael León Torres}
\tagline{Ingeniero de Software Senior, especializado en .NET, Angular y SQL Server}
\socialinfo{%
  \smartphone{+52 551 826 3805}
  \email{r.leontorres1988@gmail.com}
  \address{Salamanca, Guanajuato}\\
  \linkedin{rafael-leon-torres-75410999}
  \github{RafaelLeonTorres}\\
  \website{https://rafaelleontorres.github.io/mi-sitio2/}{rafaelleontorres.github.io/mi-sitio2}%
}


\begin{document}
\raggedright
\makecvheader

\hfill Salamanca, Guanajuato, a \today

\conempresa{Estimado equipo de reclutamiento de \laempresa:}{Estimado equipo de reclutamiento:}

Soy ingeniero de software con más de 8 años de experiencia desarrollando y optimizando sistemas críticos con C\#/.NET, Angular y SQL Server en salud, manufactura, banca y facturación electrónica. Me interesa sumarme a \conempresa{\laempresa}{su equipo} \conpuesto{en el puesto de \elpuesto}{en un puesto de desarrollo de software}, donde puedo aportar desde el análisis de los requerimientos hasta la puesta en producción.

Actualmente, en Qualis Technologies, trabajo en EMA Salud, una plataforma para clínicas, hospitales y farmacias que usan instituciones públicas de salud en México. Desarrollo funcionalidades de punta a punta con .NET 9 y Angular 19, optimizo consultas y reportes críticos sin alterar sus resultados, escribo scripts de migración de datos que se pueden volver a ejecutar sin riesgo y resuelvo incidencias productivas con análisis de causa raíz.

Antes, en NTT DATA, dirigí al equipo que migró las aplicaciones de escritorio y las APIs de Citibanamex a una nueva infraestructura durante la separación del banco. En CECOBAN participé en la migración a CFDI 4.0 y en el complemento Carta Porte 2.0, y en Qualtop desarrollé APIs .NET y aplicaciones Angular para Jabil, con integración y despliegue continuo en Azure DevOps.

Me importa que el código se pueda mantener y revisar: aplico SOLID y patrones de diseño, escribo pruebas, reviso pull requests con criterios explícitos y documento para que cualquiera pueda continuar el trabajo. También uso IA (Claude Code y servidores MCP) para acelerar el análisis y las revisiones, y la validación final siempre es mía. Explico los problemas técnicos en lenguaje claro a negocio y a QA.

Me encantaría platicar sobre cómo puedo contribuir a \conempresa{\laempresa}{su equipo}. En mi sitio encontrará mi CV y algunas publicaciones técnicas: \href{https://rafaelleontorres.github.io/mi-sitio2/}{rafaelleontorres.github.io/mi-sitio2}. Gracias por su tiempo y consideración.

Atentamente,

\medskip
\textbf{Rafael León Torres}

\end{document}
